\documentclass[11pt]{article}
\usepackage{mathblatt}
\begin{document}
% Kopfzeile Seite 1 ohne Kennung; die Rückseite (Lösungen, für den Lehrer) bekommt sie nach \begleitteil
\blattkopf{Mathematik}{Basisaufgaben}
\einheitenkopf[][Zettel 12]{Basisaufgaben · Zettel 12}
\zweigzeile{je Aufgabe 1 Punkt · Lösungen auf der Rückseite}
\par\vspace{2pt}\noindent Name \underline{\hspace{7cm}}\hfill Datum \underline{\hspace{3cm}}\hfill Punkte \underline{\hspace{1.2cm}} / 18\par\vspace{6pt}
% \small: volle Seite, kurze Aufgaben zweispaltig (v0.6)
\small

\noindent\begin{minipage}[t]{0.485\linewidth}\raggedright
% 1: Zahlen in verschiedenen Darstellungen vergleichen – brueche-dezimalzahlen-basis-k4-v10 (Original 2023-OS-B1f, ausgedacht: keine Marke)
\begin{aufgabe}{Welcher Wert ist der kleinste: $1{,}5$; $150\,\%$; $1{,}2^2$; $1{,}4$? \leerfeld}
\end{aufgabe}
\end{minipage}\hfill\begin{minipage}[t]{0.485\linewidth}\raggedright
% 2: Termwert berechnen – rationale-zahlen-basis-k1-v8 (Original 2026-FOR-B1g, ausgedacht: keine Marke)
\begin{aufgabe}{Berechne den Wert von $6 \cdot (x + 1)$ für $x = -4$. \leerfeld}
\end{aufgabe}
\end{minipage}\par

\noindent\begin{minipage}[t]{0.485\linewidth}\raggedright
% 3: Zehnerpotenzschreibweise umwandeln – potenzen-wurzeln-basis-k3-v5 (Original 2025-OS-B1d, ausgedacht: keine Marke)
\begin{aufgabe}{$27\,000\,000 = 2{,}7 \cdot 10^{\square}$ – welche Hochzahl? \leerfeld}
\end{aufgabe}
\end{minipage}\hfill\begin{minipage}[t]{0.485\linewidth}\raggedright
% 4: Prozentwert berechnen – prozentrechnung-basis-k4-v8 (Original 2026-FOR-B1a, ausgedacht: keine Marke)
\begin{aufgabe}{$5\,\%$ von $140\,€$ \leerfeld[€]}
\end{aufgabe}
\end{minipage}\par

\noindent\begin{minipage}[t]{0.485\linewidth}\raggedright
% 5: Zeiteinheiten umrechnen – einheiten-basis-k5-v1 (Original 2025-OS-B1b, ausgedacht: keine Marke)
\begin{aufgabe}{Eine Busfahrt dauert $2{,}5$ h. Gib die Dauer in Minuten an. \leerfeld[min]}
\end{aufgabe}
\end{minipage}\hfill\begin{minipage}[t]{0.485\linewidth}\raggedright
% 6: Proportionale Zuordnung Dreisatz – zuordnungen-basis-k2-v1 (Original 2023-OS-B1a, ausgedacht: keine Marke)
\begin{aufgabe}{$4$ Brötchen kosten $1{,}60$ €. Wie viel kosten $7$ Brötchen? \leerfeld[€]}
\end{aufgabe}
\end{minipage}\par

\noindent\begin{minipage}[t]{0.485\linewidth}\raggedright
% 7: Lineare Gleichung lösen – lineare-gleichungen-basis-k1-v7 (Original 2024-OS-B1d, ausgedacht: keine Marke)
\begin{aufgabe}{Löse die Gleichung $2 \cdot (x + 4) - 3 = 5$. x = \leerfeld}
\end{aufgabe}
\end{minipage}\hfill\begin{minipage}[t]{0.485\linewidth}\raggedright
% 8: Quadratseite aus Fläche berechnen – flaechen-basis-k2-v1 (Original 2020-OS-B1d, ausgedacht: keine Marke)
\begin{aufgabe}{Eine quadratische Fliese hat den Flächeninhalt $49$ cm². Wie lang ist eine Seite? \leerfeld[cm]}
\end{aufgabe}
\end{minipage}\par

\noindent\begin{minipage}[t]{0.485\linewidth}\raggedright
% 9: Median bestimmen – daten-basis-k3-v1 (Original 2024-OS-B1h, ausgedacht: keine Marke)
\begin{aufgabe}{Tore in fünf Spielen: $3$; $1$; $4$; $0$; $2$. Bestimme den Median. \leerfeld}
\end{aufgabe}
\end{minipage}\hfill\begin{minipage}[t]{0.485\linewidth}\raggedright
% 10: Gleichschenkliges Dreieck erkennen – winkel-dreiecke-basis-k2-v2 (Original 2023-OS-B1g, ausgedacht: keine Marke)
\begin{aufgabe}{Im Dreieck ABC ist $\alpha = \beta = 75^\circ$ und AC = $8$ cm. Gib die Länge von BC und die Größe von $\gamma$ an. BC = \leerfeld[cm] , $\gamma =$ \leerfeld °}
\end{aufgabe}
\end{minipage}\par

% 11: Prozent und Anteil umwandeln – prozentrechnung-basis-k2-v3 (Original 2022-OS-B1f, ausgedacht: keine Marke)
\begin{aufgabe}{3 von 4 Plätzen sind belegt – wie viel Prozent? Kreuze an. \\ \kreuz{$3\,\%$} \kreuz{$34\,\%$} \kreuz{$60\,\%$} \kreuz{$75\,\%$}}
\end{aufgabe}

% 12: Term zu Sachtext angeben – terme-basis-k2-v6 (Original 2023-OS-B1h, ausgedacht: keine Marke)
\begin{aufgabe}{Das Dreifache einer Zahl $x$ wird von 20 subtrahiert. Welcher Term passt? \\ \kreuz{$3x - 20$} \kreuz{$20 - 3x$} \kreuz{$3 \cdot (20 - x)$} \kreuz{$20x - 3$}}
\end{aufgabe}

% 13: Lage eines Punktes zu den Achsen erkennen – symmetrie-abbildungen-basis-k2-v2 (Original 2020-OS-B1b, ausgedacht: keine Marke)
\begin{aufgabe}{Genau einer der vier Punkte liegt auf der y-Achse. Welcher? \\ \kreuz{$A(2|6)$} \kreuz{$B(0|6)$} \kreuz{$C(6|0)$} \kreuz{$D(-2|-6)$}}
\end{aufgabe}

% 14: Kantenzahl eines Körpers angeben – koerper-basis-k2-v2 (Original 2019-OS-B1f, ausgedacht: keine Marke)
\begin{aufgabe}{Eine Schokoladenpackung hat die Form eines Prismas mit dreieckiger Grundfläche. Gib an, wie viele Kanten sie hat. \leerfeld Kanten}
\end{aufgabe}

% 15: Wahrscheinlichkeit einstufig – bruchrechnung-basis-k1-v4 (Original 2016-OS-B1f, ausgedacht: keine Marke)
\begin{aufgabe}{In einem Beutel liegen 3 rote, 4 blaue und 5 grüne Kugeln. Eine Kugel wird gezogen. Wie groß ist die Wahrscheinlichkeit, dass sie nicht grün ist? \leerfeld}
\end{aufgabe}

% 16: Zufallsgerät zu Wahrscheinlichkeit entwerfen – wahrscheinlichkeit-basis-k2-v1 (Original 2019-OS-B1g, ausgedacht: keine Marke)
\begin{aufgabe}{Ein Glücksrad hat $8$ gleich große Felder. Die Wahrscheinlichkeit für Blau soll $25\,\%$ sein. Wie viele Felder müssen blau sein? \leerfeld Felder}
\end{aufgabe}

% 17: Körper aus Netz oder Schrägbild benennen – koerper-basis-k4-v1 (Original 2018-OS-B1i, ausgedacht: keine Marke)
\begin{aufgabe}{\begin{minipage}[t]{0.6\linewidth}Zu welchem Körper gehört das Netz? Kreuze an. \\ \kreuz{Würfel} \kreuz{Quader} \kreuz{Prisma mit dreieckiger Grundfläche}\end{minipage}\hfill\begin{minipage}[t]{0.37\linewidth}\vspace{0pt}
\netzquader{2}{1}{1.5}{}{}{}
\end{minipage}}
\end{aufgabe}

% 18: Pythagoras Gleichung zuordnen – pythagoras-basis-k1-v1 (Original 2026-FOR-B1j, ausgedacht: keine Marke)
\begin{aufgabe}{\begin{minipage}[t]{0.6\linewidth}Welche Gleichung gilt in diesem rechtwinkligen Dreieck? \\ \kreuz{$r^2 = p^2 + q^2$} \kreuz{$p^2 = q^2 + r^2$} \kreuz{$r^2 = p^2 - q^2$} \kreuz{$r = p + q$}\end{minipage}\hfill\begin{minipage}[t]{0.37\linewidth}\vspace{0pt}
\dreieckrw{3.5}{2.5}{p}{q}{r}
\end{minipage}}
\end{aufgabe}

\begleitteil
\blattkopf{Basisaufgaben · Zettel 12}{BAS-Z12 · Lösungen}
\erg{1}{$1{,}4$; die anderen sind $1{,}5$, $1{,}5$ und $1{,}2^2 = 1{,}44$ \hfill {\footnotesize\textit{Zahlen in verschiedenen Darstellungen vergleichen · 2023}}}
\erg{2}{$6 \cdot (-4 + 1) = 6 \cdot (-3) = -18$ \hfill {\footnotesize\textit{Termwert berechnen · 2026}}}
\erg{3}{$7$ – das Komma wandert sieben Stellen nach links \hfill {\footnotesize\textit{Zehnerpotenzschreibweise umwandeln · 2025}}}
\erg{4}{$0{,}05 \cdot 140 = 7\,€$ \hfill {\footnotesize\textit{Prozentwert berechnen · 2026}}}
\erg{5}{$2{,}5 \cdot 60 = 150$ min \hfill {\footnotesize\textit{Zeiteinheiten umrechnen · 2025}}}
\erg{6}{Für $1$: $1{,}60 : 4 = 0{,}40$; $0{,}40 \cdot 7 = 2{,}80\,€$ \hfill {\footnotesize\textit{Proportionale Zuordnung Dreisatz · 2023}}}
\erg{7}{$2 \cdot (x + 4) = 8$, $x + 4 = 4$, also $x = 0$ \hfill {\footnotesize\textit{Lineare Gleichung lösen · 2024}}}
\erg{8}{$7 \cdot 7 = 49$, also $7$ cm \hfill {\footnotesize\textit{Quadratseite aus Fläche berechnen · 2020}}}
\erg{9}{geordnet $0$; $1$; $2$; $3$; $4$; der mittlere Wert ist $2$ \hfill {\footnotesize\textit{Median bestimmen · 2024}}}
\erg{10}{BC = AC = $8$ cm (gleiche Winkel bei A und B); $\gamma = 180^\circ - 2 \cdot 75^\circ = 30^\circ$ \hfill {\footnotesize\textit{Gleichschenkliges Dreieck erkennen · 2023}}}
\erg{11}{$75\,\%$ ($\frac{3}{4} = \frac{75}{100}$) \hfill {\footnotesize\textit{Prozent und Anteil umwandeln · 2022}}}
\erg{12}{$20 - 3x$ \hfill {\footnotesize\textit{Term zu Sachtext angeben · 2023}}}
\erg{13}{$B(0|6)$ \hfill {\footnotesize\textit{Lage eines Punktes zu den Achsen erkennen · 2020}}}
\erg{14}{$3 + 3 + 3 = 9$ Kanten ($3$ oben, $3$ unten, $3$ Seitenkanten) \hfill {\footnotesize\textit{Kantenzahl eines Körpers angeben · 2019}}}
\erg{15}{$\frac{7}{12}$ (7 von 12 Kugeln sind rot oder blau) \hfill {\footnotesize\textit{Wahrscheinlichkeit einstufig · 2016}}}
\erg{16}{$25\,\% = \frac{1}{4}$; $\frac{1}{4} \cdot 8 = 2$ \hfill {\footnotesize\textit{Zufallsgerät zu Wahrscheinlichkeit entwerfen · 2019}}}
\erg{17}{Quader \hfill {\footnotesize\textit{Körper aus Netz oder Schrägbild benennen · 2018}}}
\erg{18}{$r^2 = p^2 + q^2$ \hfill {\footnotesize\textit{Pythagoras Gleichung zuordnen · 2026}}}
\end{document}
